\section{Joint cuts in the original variables}\label{sec:aggregation}

Let the original model contain the inequality sides
\begin{equation}\label{eq:sides}
g_i(x)+\ell_i(v)\le b_i,\qquad i=1,\ldots,m,
\end{equation}
where $x$ is a selected subset of original variables, $v$ contains model
variables, and each $\ell_i$ is linear. A variable may occur both in $x$
and in an affine remainder. All constants are included in $b_i$.
The block domain $D$ contains the block coordinates of every feasible
point to which a proposed cut will apply.

\begin{proposition}[Original-variable support cut]\label{prop:aggregate}
Choose $a\in\R^{\dim x}$ and $\lambda\in\R^m_+$. If
\[
\beta\le\inf_{u\in D}\left\{a^\top u+
                         \sum_i\lambda_i g_i(u)\right\},
\]
then the original-variable linear inequality
\begin{equation}\label{eq:aggregate}
a^\top x-\sum_i\lambda_i\ell_i(v)
\ge\beta-\sum_i\lambda_i b_i
\end{equation}
is valid for the model.
\end{proposition}
\begin{proof}
At a feasible point, nonnegativity of the multipliers gives
$\sum_i\lambda_i g_i(x)\le\sum_i\lambda_i(b_i-\ell_i(v))$.
Substitute this upper bound into the support inequality and rearrange.
\end{proof}

This is weak Lagrangian duality. It needs neither convex functions nor
optimal multipliers. The functions are minimized jointly before elimination;
simply combining already linear valid inequalities cannot strengthen their
intersection. Equality sides can use unrestricted multipliers, or two
opposite inequality sides with nonnegative multipliers. A minimization
objective $g_0(x)+\ell_0(v)+c_0$ supplies an epigraph side
$g_0(x)+\ell_0(v)-t\le-c_0$. A maximization objective uses the
hypograph side $-g_0(x)-\ell_0(v)+t\le c_0$. The objective
transformation, constant, and sense must be identical in every compared
mode. These rows do not require a separate equality variable for each
selected nonlinear expression.

\begin{theorem}[Closure of the direct-row interface]\label{thm:closure}
Suppose $D$ is nonempty and compact and $g:D\to\R^m$ is continuous. Put
\[
K=\conv\{(x,g(x)):x\in D\},\qquad
P=\{(x,z):\exists y,\ (x,y)\in K,\ y+Lz\le b\}.
\]
Then $P$ is exactly the intersection, over all $a\in\R^{\dim x}$ and
$\lambda\in\R^m_+$, of
\begin{equation}\label{eq:closure}
a^\top x-\lambda^\top Lz\ge
\min_{u\in D}\{a^\top u+\lambda^\top g(u)\}-\lambda^\top b.
\end{equation}
If continuous equality features $h:D\to\R^r$ are included in the
joint graph, their support coefficients are free and their target values
are substituted without adding cone coordinates.
\end{theorem}
\begin{proof}
The set $U=K+(\{0\}\times\R^m_+)$ is closed and convex. To see
closedness, take a convergent sequence of sums, extract a convergent
subsequence of the compact $K$ summands, and subtract; the cone summands
then converge inside their closed cone. Membership in $P$ is equivalent
to $(x,b-Lz)\in U$.

A lower support functional of $U$ with finite constant must have
nonnegative coefficients on the last $m$ coordinates, because a negative
coefficient has infimum $-\infty$ along the corresponding cone ray.
For the remaining directions its support value is
$\min_{u\in D}(a^\top u+\lambda^\top g(u))$. Substituting
$(x,b-Lz)$ gives~\eqref{eq:closure}. Closed convex separation excludes
every point outside $U$ by such a finite support inequality, proving the
converse. For continuous equality features the joint graph
$\conv\{(x,g(x),h(x)):x\in D\}$ is also compact. The same argument
adds no cone in their coordinates, so their coefficients are free.
\end{proof}

This describes the projection of the specified graph relaxation, not
necessarily the convex hull of original feasible points. For example,
on $[0,1]$ the equality $x^2=1/4$ has feasible set $\{1/2\}$.
The mixture with mass $3/4$ at zero and $1/4$ at one instead has
$\E x=\E x^2=1/4$. The projected graph relaxation thus contains
$x=1/4$, and every direct support cut accepts it. Including $x^2=1/4$
inside the support domain changes the oracle problem and removes this
example. Convexifying a graph and imposing constraints in expectation do
not generally commute. Ordinary domain propagation can also resolve this
particular equality; that does not change the comparison of cut families
over the unchanged box.

\begin{proposition}[Safe export after elimination]\label{prop:elimination-round}
Let the exactly eliminated row be $c^\top v\ge r$ and let
$\widehat c$ denote the exact rational values of the coefficients
actually exported in binary64. With $\Delta_j=\widehat c_j-c_j$, suppose
valid global bounds imply a finite value
\[
E=\sum_j\inf_{L_j\le s\le U_j}\Delta_j s.
\]
Then every exported lower side $\widehat r\le r+E$ gives a valid row
$\widehat c^\top v\ge\widehat r$.
\end{proposition}
\begin{proof}
$\widehat c^\top v=c^\top v+\sum_j\Delta_jv_j\ge r+E$.
\end{proof}

A positive $\Delta_j$ needs a finite lower bound, a negative one a
finite upper bound, and a zero one contributes zero even on an unbounded
coordinate. This covers an unbounded objective epigraph when its exported
coefficient is exact. The correction must follow coefficient merging,
normalization, and any zero dropping. The lower side is rounded downward
and checked by exact rational comparison. An unavailable finite correction,
overflow, or lost useful separation causes refusal; a fixed safety shift
cannot replace the calculation. The source sides, their orientations,
multipliers, domain provenance, support bound, exact eliminated row, and
actual exported row are separate parts of the certificate contract.

The derivation and its checked implementation are detailed in
\pathref{../theory/row-aggregation.md}{theory/row-aggregation.md}.
The formulation and safe-rounding principles have direct prior literature;
Section~\ref{sec:literature} records the sources and contribution boundary.
